\section{Nguyên lý các thuật toán}
Trong đồ án, nhóm em cài đặt và so sánh các thuật toán sau. Bảng dưới đây tóm tắt điều kiện áp dụng trên graph hữu hạn và cost không âm; riêng các thuật toán tour được đánh giá với số điểm giao hàng nhỏ.

{\footnotesize
\begin{tabularx}{\textwidth}{>{\raggedright\arraybackslash}p{2.0cm} >{\raggedright\arraybackslash}X >{\raggedright\arraybackslash}p{2.4cm} >{\raggedright\arraybackslash}p{1.2cm} >{\raggedright\arraybackslash}p{1.2cm} >{\raggedright\arraybackslash}p{1.2cm}}
\toprule
Thuật toán & Quy tắc chọn node & Thời gian & Bộ nhớ & Đầy đủ? & Tối ưu?\\
\midrule
BFS & Mở rộng theo từng lớp & \(O(V+E)\) & \(O(V)\) & Có & Có\\
DFS & Đi sâu một nhánh trước & \(O(V+E)\) & \(O(V)\) & Có & Không\\
UCS/Dijkstra & Chọn \(g\) nhỏ nhất & \(O((V+E)\log V)\) & \(O(V)\) & Có & Có\\
A* & Chọn \(f=g+h\) nhỏ nhất & Phụ thuộc \(h\) & \(O(V)\) & Có & Có\\
Greedy Best-First & Chọn \(h\) nhỏ nhất & Phụ thuộc \(h\) & \(O(V)\) & Có & Không\\
Bidirectional & Tìm từ hai đầu & Thường giảm & \(O(b^{d/2})\) & Có & Có\\
\bottomrule
\end{tabularx}
}

Các kết luận trong bảng được hiểu với graph hữu hạn và cách cài đặt có kiểm tra trạng thái đã thăm. BFS tối ưu khi mọi edge có cùng cost, nên trong bài này chủ yếu dùng để so sánh. DFS có thể tìm được đường trên graph hữu hạn nhưng không đảm bảo đó là đường rẻ nhất. UCS được nhóm em dùng làm mốc tham chiếu vì tối ưu với mọi edge cost không âm. A* có cùng bảo đảm với UCS khi heuristic admissible và consistent; đây là lý do nhóm em kiểm tra A* với UCS thay vì chỉ nhìn thời gian chạy. Greedy thường nhanh nhưng không thể dùng để kết luận nghiệm tối ưu. Bidirectional cần điều kiện dừng đúng và cách xử lý graph có hướng rõ ràng.

Với bài toán đi qua nhiều điểm, Held--Karp là dynamic programming chính xác với độ phức tạp \(O(n^2 2^n)\), còn Nearest Neighbor là heuristic tham lam khoảng \(O(n^2)\) và không bảo đảm tối ưu. Vì vậy nhóm em dùng Held--Karp làm mốc khi số stop nhỏ và dùng Nearest Neighbor để minh hoạ sự đánh đổi giữa chất lượng tour và thời gian chạy.

\begin{center}
{\scriptsize
\begin{tabular}{l l c c}
\toprule
Thuật toán tour & Cách chọn & Thời gian & Tối ưu?\\
\midrule
Held--Karp & Lưu nghiệm tốt nhất theo tập stop đã đi qua & \(O(n^2 2^n)\) & Có\\
Nearest Neighbor & Chọn stop tiếp theo có pairwise cost nhỏ nhất & \(O(n^2)\) & Không\\
\bottomrule
\end{tabular}
}
\end{center}

Trong report, A* được kết luận là tối ưu vì dùng lower-bound heuristic của Chương 6. Nếu thay bằng heuristic tùy ý thì cần kiểm tra lại tính admissible và consistent. Bidirectional cũng chỉ tối ưu khi hai hướng dùng cost tương thích và điều kiện dừng đúng.

\begin{figure}[htbp]
\centering
\begin{tikzpicture}[>=stealth, node distance=17mm, every node/.style={font=\small}, point/.style={circle, draw, fill=blue!8, minimum size=8mm}]
\node[point] (n1) {N01};
\node[point, right=of n1] (n2) {N02};
\node[point, right=of n2] (n6) {N06};
\node[point, below=of n2] (n4) {N04};
\draw[->] (n1) -- node[above] {E01} (n2);
\draw[->] (n2) -- node[above] {E03} (n6);
\draw[->] (n2) -- node[left] {E11} (n4);
\draw[->] (n4) -- node[below] {E09} (n6);
\end{tikzpicture}
\caption{Toy graph minh hoạ hai phương án từ N01 đến N06. Thuật toán tìm kiếm đánh giá route theo tổng composite cost, không chỉ theo số cạnh.}
\label{fig:toy-search-graph}
\end{figure}
